\chapter*{Resumen}
\addcontentsline{toc}{chapter}{Resumen}

\noindent
Esta tesis desarrolla los modelos de orden reducido intrusivos basados en proyección (PROM) como base escalable y con consistencia física para habilitar gemelos digitales industriales con operación en tiempo real. Al integrar métodos numéricos de alta fidelidad con técnicas sólidas de reducción de modelos, aborda los cuellos de botella computacionales que dificultan el despliegue de simulaciones a gran escala para diseño, optimización y monitorización de la operación.

Estructurada como un compendio de artículos revisados por pares, la tesis presenta tres contribuciones metodológicas principales. En primer lugar, se propone un marco de hiperreducción para PROM de Petrov-Galerkin que elimina la necesidad de mallas complementarias y permite un muestreo a nivel de elemento plenamente compatible con los flujos de trabajo estándar de elementos finitos (por ejemplo, Kratos Multiphysics). En segundo lugar, se plantea un flujo de trabajo escalable habilitado para computación de alto rendimiento (HPC) para el entrenamiento y el despliegue de PROM, que aprovecha la generación paralela de instantáneas (snapshots), la descomposición en valores singulares (SVD) distribuida y una versión paralela del método de cubatura empírica (ECM), demostrada en un gemelo digital térmico industrial de un motor eléctrico. En tercer lugar, se extienden las metodologías PROM a regímenes no lineales mediante estrategias de cierre en espacio latente, incluidas PROM-ANN (redes neuronales artificiales) y modelos sustitutos interpretables basados en kernel, como PROM-GPR (regresión por procesos gaussianos) y PROM-RBF (funciones de base radial), con el fin de superar la barrera de la anchura n de Kolmogórov en flujos dominados por la convección. Esta extensión incorpora una estrategia de entrenamiento informada por la física discreta que alinea las variedades aprendidas con redes neuronales con el comportamiento del residuo del esquema numérico subyacente y garantiza la consistencia física.

Los métodos desarrollados se validan de forma sistemática en problemas canónicos, como la ecuación de Burgers unidimensional inviscida, y se trasladan a configuraciones industriales como la estela del cuerpo de Ahmed mediante el entorno AERO-F. Asimismo, se revisan técnicas clásicas, como las variedades lineales a trozos y las cuadráticas, para contextualizar las limitaciones de los subespacios lineales globales y motivar las estrategias no lineales propuestas.

Todas las contribuciones se han implementado en los marcos de código abierto Kratos Multiphysics y AERO-F, lo que subraya su aplicabilidad práctica en flujos de trabajo de ingeniería a gran escala. Aunque esta tesis no implementa un gemelo digital de bucle cerrado, proporciona una base sólida para futuros gemelos digitales de componente y de activo, y tiende un puente entre la modelización de alta fidelidad y las capacidades predictivas en tiempo real.

Por último, la tesis impulsa la transferencia tecnológica mediante la iniciativa SimTwins, que ofrece sombras digitales y gemelos digitales escalables y de código abierto, adaptados a las demandas emergentes de la Industria 4.0 y 5.0. En conjunto, estas contribuciones sitúan a los PROM intrusivos basados en proyección como herramientas robustas, interpretables y de alto rendimiento, esenciales para los ecosistemas de gemelos digitales de próxima generación.

